This section summarizes implementation decisions and the artifacts we provide to support reproducibility. We describe (i) model interfaces and prompt template organization, (ii) the strict JSON schema that leaf-section generators must emit, (iii) the per-node evidence-injection mechanism, (iv) document-graph construction and the scheduler interface, and (v) the LaTeX assembly and refs.bib generation pipeline. Where appropriate we emphasize design choices made to prioritize auditability and schema validation over unconstrained free-text output.

Model interfaces and prompt templates.
The implementation is organized around a small set of interoperable components with well-defined I/O contracts rather than a single monolithic program. Concretely, a generator component accepts a JSON-formatted prompt payload and returns a JSON-conforming response; a retriever component accepts a textual query and returns a ranked list of evidence items (each with a short snippet and a canonical source pointer); and an assembler component consumes validated section JSONs and emits LaTeX fragments plus a BibTeX file. Prompt templates are stored as parameterized plain-text strings that are rendered with (a) the section node metadata (heading, target length, role), (b) contextual graph neighbors (parent/child headings and short summaries when available), and (c) a numbered list of retrieved evidence items. Prompts include explicit instruction blocks that require the generator to (1) produce only the fields mandated by the JSON schema, (2) associate every non-trivial factual claim with one or more evidence item indices, and (3) avoid speculation beyond the provided evidence. Prompt files and representative prompt–response examples are persisted in the release so that template text and any subsequent edits are inspectable and diffable.

JSON schema enforced for section outputs.
We require each leaf-section generator to emit a single JSON object conforming to a strict machine-readable schema. The schema codifies fields for section metadata (id, heading, role, target length), an explicit claim inventory (each claim includes an identifier, a short natural-language statement, and one or more evidence references), paragraph text with span-level indices that may be associated to claims, a list of citation objects that map spans to evidence items, and provenance metadata that records retrieval queries and retrieved-item identifiers. Schema validation is performed automatically at generation time; outputs that fail validation are rejected and the generator is prompted for correction. We include the exact JSON Schema file(s) together with exemplar valid section JSONs in the project release so that downstream users can reproduce the validator behavior and run automated checks. Our use of an explicit, verifiable schema is motivated by the central role that JSON Schema plays as a standard for describing families of JSON documents and by the non-trivial static-analysis challenges that attend JSON Schema reasoning; prior work has explored the complexity of schema satisfiability, witness generation and related analyses for expressive JSON Schema fragments \cite{web_attouche_2022_witness_generation_for_json_schema,web_baazizi_2021_not_elimination_and_witness_generation_for_json_}.

Evidence-injection mechanism per node.
For each scheduled generation step the system constructs one or more retrieval queries from (a) the section node text (heading, role, short intent), (b) context drawn from dependent nodes as specified by the document graph, and (c) any accumulated claims in neighboring nodes when applicable. The retriever returns a ranked list of items; each item carries a condensed snippet, a canonical source pointer (an index identifier), and metadata (document title, page, paragraph). These items are injected into the prompt as a numbered list; generators refer to evidence items by number when populating claim evidence_refs fields. Prompts include explicit constraints designed to discourage hallucination (for example: "Do not assert facts about which none of the provided evidence items provide support") and require that each substantive claim include at least one evidence reference. To aid reproducibility we release a small calibration set of queries paired with expected retrieved items so that developers can verify that their local retrieval index and query construction heuristics produce compatible results.

Document-graph construction and scheduler interface.
The document graph is represented as a directed acyclic graph whose nodes denote sections and whose typed edges encode dependency relations (e.g., prerequisite, elaboration, evidence-flow). Each node carries scheduling metadata (id, heading, role, target_length) and optional short summaries to facilitate prompt construction. The scheduler exposes a programmatic API that, given a graph, returns an ordered sequence of generation batches; the API supports both single-threaded traversal and batch emission suitable for parallel generation of independent subtrees. Our reference scheduling heuristic is dependency-respecting and favors producing leaf-node outputs earlier when those outputs are required as context for parents, but the scheduler implementation is configurable and supports alternative traversal strategies (top-down, breadth-first, etc.). The scheduler source code and configuration files used for the experiments reported in this paper are included in the release to enable re-running and comparing alternative policies.

Assembly into LaTeX and refs.bib.
After section JSONs are validated, an assembler module converts structured paragraphs and citation slots into LaTeX fragments. Each citation object that contains an explicit cite_key is mapped directly to a BibTeX entry; when generators provide only a retrieval source pointer the assembler resolves or synthesizes a BibTeX entry using retriever-provided metadata and records the mapping in refs.bib. Bibliographic deduplication is performed conservatively by normalizing titles and author lists and flagging borderline matches for manual adjudication. Cross-reference links (for example, references to other sections by node id) are translated into \verb|\label| and \verb|\ref| pairs. The assembler additionally emits a machine-readable audit file that pairs every claim id with its evidence_refs and the corresponding source pointers; this audit trail enables both automated verification and human review of claim–evidence alignment.

Reproducibility artifacts and expected contents.
To facilitate reproduction we release the following artifacts alongside this paper: (1) the JSON Schema files and exemplar valid section JSONs; (2) prompt templates and a small set of prompt–response example pairs used during development; (3) the scheduler source and configuration files used in our reported runs; (4) a minimal retrieval index and scripts to build it from the curated PDF set described in Section 4.1 (where licensing permits); and (5) the assembler code that maps validated JSONs to LaTeX and generates refs.bib. Where licensing or privacy constraints prevent redistribution of some internal retrieval sources, we provide synthetic stand-ins and scripts that demonstrate how to construct a compatible index. Statements about artifact contents reflect the planned contents of our release; users should consult the project repository for the actual files accompanying this submission.

Compute and runtime considerations (general background knowledge).
Reproducing generation experiments requires (at minimum) access to a language model and a retrieval backend; our implementation is model-agnostic and separates the generator API from the concrete model provider so practitioners may substitute smaller or open models for experimentation (general background knowledge). Wall-clock runtime depends primarily on (a) the number of sections generated, (b) average retrieval latency per node, and (c) the degree of parallelism configured for batch generation; to support trade-offs we provide configuration knobs (batch sizes, concurrency limits, retrieval cache settings) and guidance scripts that tune for either lower latency or higher fidelity (general background knowledge).

Limitations and notes.
Because some external retrieval or runtime artifacts cannot be embedded directly in the paper, environment variables, API keys, and dataset access permissions are documented in the artifact README rather than in the paper body. Reproduction across different underlying language models may yield variation in surface phrasing and in the frequency of evidence-linked claims; our schema validation and audit trail are intended to mitigate such variability by making outputs machine-verifiable. Finally, we note that the statements about JSON Schema complexity and the motivation for strict schema validation are supported by prior analyses of JSON Schema reasoning and witness generation \cite{web_attouche_2022_witness_generation_for_json_schema,web_baazizi_2021_not_elimination_and_witness_generation_for_json_}.